\subsection{Exact sequences}

Recall the following definition (Algebra, Chapter II, § 1, no. 4):

DEFINITION 1. Let A be a ring, F, G, H three right (resp. left) A-modules, f a homomorphism of F to G and g a homomorphism of G to H. The ordered pair \((f, g)\) is called an exact sequence if \(\bar{g}^{-1}(0) = f(F)\) , that is if the kernel of g is equal to the image of f.

The diagram

\[
\mathrm{F} \xrightarrow {f} \mathrm{G} \xrightarrow {g} \mathrm{H}\tag{4}
\]

is also called an exact sequence.

Consider similarly a diagram consisting of four modules and three homomorphisms:

\[
\mathrm{E} \xrightarrow {f} \mathrm{F} \xrightarrow {g} \mathrm{G} \xrightarrow {h} \mathrm{H}.\tag{5}
\]

This diagram is said to be exact at F if the diagram \(E \xrightarrow{f} F \xrightarrow{g} G\) is an exact sequence; it is said to be exact at G if \(F \xrightarrow{g} G \xrightarrow{h} H\) is an exact sequence. If (5) is exact at F and G, it is said to be exact or also an exact sequence. Exact sequences with any number of terms are similarly defined.

Recall the following results (loc. cit.), where E, F, G denote right (resp. left)

A-modules, the arrows represent homomorphisms and 0 denotes a module reduced to its identity element:

\begin{enumerate}
\def\labelenumi{(\alph{enumi})}
\item
  To say that \(0 \to \mathbf{E} \xrightarrow{f} \mathbf{F}\) is an exact sequence is equivalent to saying that \(f\) is injective.
\item
  To say that \(\mathbf{E} \xrightarrow{f} \mathbf{F} \to 0\) is an exact sequence is equivalent to saying that \(f\) is surjective.
\item
  To say that \(0 \to \mathbf{E} \xrightarrow{f} \mathbf{F} \to 0\) is an exact sequence is equivalent to saying that \(f\) is bijective, that is that \(f\) is an isomorphism of \(\mathbf{E}\) onto \(\mathbf{F}\) .
\item
  If F is a submodule of E and i denotes the canonical injection of F into E and p the canonical surjection of E onto E/F, the diagram
\end{enumerate}

\[
0 \longrightarrow \mathrm{F} \stackrel {i} {\longrightarrow} \mathrm{E} \stackrel {p} {\longrightarrow} \mathrm{E} / \mathrm{F} \longrightarrow 0\tag{6}
\]

is an exact sequence.

\begin{enumerate}
\def\labelenumi{(\alph{enumi})}
\setcounter{enumi}{4}
\tightlist
\item
  If \(f: \mathbf{E} \to \mathbf{F}\) is a homomorphism, the diagram
\end{enumerate}

\[
0 \longrightarrow \overline {{f}} ^ {- 1} (0) \stackrel {i} {\longrightarrow} \mathrm{E} \stackrel {f} {\longrightarrow} \mathrm{F} \stackrel {p} {\longrightarrow} \mathrm{F} / f (\mathrm{E}) \longrightarrow 0\tag{7}
\]

(where \(i\) is the canonical injection of \(\overline{f}(0)\) into E and \(p\) the canonical projection of \(\mathbf{F} / f(\mathbf{E})\) ) is an exact sequence.

\begin{enumerate}
\def\labelenumi{(\alph{enumi})}
\setcounter{enumi}{5}
\tightlist
\item
  For the diagram
\end{enumerate}

\[
\mathrm{E} \xrightarrow {f} \mathrm{F} \xrightarrow {g} \mathrm{G}\tag{8}
\]

to be an exact sequence, it is necessary and sufficient that there exist modules S, T and homomorphisms \(a: \mathrm{E} \to \mathrm{S}\) , \(b: \mathrm{S} \to \mathrm{F}\) , \(c: \mathrm{F} \to \mathrm{T}\) and \(d: \mathrm{T} \to \mathrm{G}\) such that \(f = b \circ a\) , \(g = d \circ c\) and the three sequences

\[
\begin{array}{c} \mathrm{E} \xrightarrow {a} \mathrm{S} \longrightarrow 0 \\ 0 \longrightarrow \mathrm{S} \xrightarrow {b} \mathrm{F} \xrightarrow {c} \mathrm{T} \longrightarrow 0 \\ 0 \longrightarrow \mathrm{T} \xrightarrow {d} \mathrm{G} \end{array}\tag{9}
\]

be exact.

Recall finally that if \(f: \mathbf{E} \to \mathbf{F}\) is an A-module homomorphism, we set \(\operatorname{Ker}(f) = \bar{f}^{-1}(0)\) , \(\operatorname{Im}(f) = f(\mathbf{E})\) , \(\operatorname{Coim}(f) = \mathbf{E} / \bar{f}^{-1}(0)\) and \(\operatorname{Coker}(f) = \mathbf{F} / f(\mathbf{E})\) . With this notation it is possible to take, in (9), \(\mathbf{S} = \operatorname{Im}(f) = \operatorname{Ker}(g)\) and \(\mathbf{T} = \operatorname{Im}(g)\) (canonically isomorphic to \(\operatorname{Coker}(f)\) ).

